\chapter{Causality and the Experimental Ideal}\label{ch:b2:causality-and-experiments}

Regression measures associations. An effect is a different object: what would have happened to the same people under a
different choice. This chapter defines effects in terms of outcomes that cannot all be observed, shows why comparing those
who were treated with those who were not usually measures something else, and shows what assigning the treatment at
random buys. The chart shows the result of one such experiment, a training programme for men with weak work histories.

\begin{openingchart}{Mean earnings in 1978 of men assigned at random to a training programme or to a control group, with
  95 per cent confidence intervals, dollars.}{National Supported Work sample of LaLonde (1986), subsample of Dehejia and
  Wahba (1999); our computation}
\begin{tikzpicture}
\begin{axis}[omrecord, width=0.5\linewidth, height=0.4\linewidth, ybar, bar width=28pt, xmin=-0.6, xmax=1.6, ymin=0,
    ymax=8500, xtick={0,1}, xticklabels={control (260 men), trained (185 men)}, ylabel={Dollars},
    yticklabel style={/pgf/number format/1000 sep={\,}}, nodes near coords, nodes near coords align={center},
    nodes near coords style={font=\scriptsize, anchor=south west, xshift=2pt, /pgf/number format/1000 sep={\,}}]
  \addplot[fill=omDef!70, draw=omDef, error bars/.cd, y dir=both, y explicit]
    table[x=group, y=mean, y error plus expr=\thisrow{hi}-\thisrow{mean}, y error minus expr=\thisrow{mean}-\thisrow{lo}]
    {figdata/bachelor-2/33-causality-and-experiments/out/means.csv};
\end{axis}
\end{tikzpicture}
\end{openingchart}

% figdata: means (control 4555 [3888, 5221], treated 6349 [5215, 7483])
The trained men earned 6,349 dollars on average in 1978 and the control men 4,555: a difference of 1,794 dollars. Because
a lottery decided who was trained, the two groups were alike in everything but the training, and the difference measures
its effect. The same trained men compared with men from a general survey would seem to have lost 8,498 dollars.

\begin{recall}
From the Year 1 volume: random samples, the difference of two means and its standard error, the omitted variable. From
\cref{ch:b2:multiple-regression}: the \omterm{def:b2:multiple-regression:ovb}{omitted-variable bias} formula. From \cref{ch:b2:regression-inference}: robust
standard errors.
\end{recall}

\section{Potential outcomes}\label{sec:b2:causality-and-experiments:potential}

\begin{definition}[Potential outcome]\label{def:b2:causality-and-experiments:potential-outcome}
The \emph{potential outcomes}\index{potential outcome} of unit $i$ are $Y_i(1)$, the outcome it would have if treated,
and $Y_i(0)$, the outcome it would have if not. Only one of them is observed: $Y_i = D_iY_i(1) + (1 - D_i)Y_i(0)$, where
$D_i$ is 1 if $i$ is treated.
\end{definition}

\begin{definition}[Treatment effect]\label{def:b2:causality-and-experiments:treatment-effect}
The \emph{treatment effect}\index{treatment effect} on unit $i$ is $Y_i(1) - Y_i(0)$. It is never observed, since no unit
is both treated and untreated at the same time.
\end{definition}

\begin{definition}[Stable unit treatment value assumption]\label{def:b2:causality-and-experiments:sutva}
The \emph{stable unit treatment value assumption}\index{stable unit treatment value assumption} states that each unit's
\omterm{def:b2:causality-and-experiments:potential-outcome}{potential outcomes} depend only on its own treatment, not on who else is treated, and that there is one version of the
treatment.
\end{definition}

The assumption fails when treated workers compete with untreated ones for the same jobs, or when vaccinated people protect
unvaccinated ones; the effects defined above then depend on how many are treated.

\section{Average effects}\label{sec:b2:causality-and-experiments:averages}

\begin{definition}[Average treatment effect]\label{def:b2:causality-and-experiments:ate}
The \emph{average treatment effect}\index{average treatment effect} is $\mathit{ATE} = E[Y_i(1) - Y_i(0)]$, the mean
effect over the whole population.
\end{definition}

\begin{definition}[Average treatment effect on the treated]\label{def:b2:causality-and-experiments:att}
The \emph{average treatment effect on the treated}\index{average treatment effect on the treated} is $\mathit{ATT} =
E[Y_i(1) - Y_i(0) \mid D_i = 1]$, the mean effect among those who were treated.
\end{definition}

The two differ when effects differ across people and the treated are not a random part of the population: if those who
gain most from training are the ones who take it, the ATT exceeds the ATE. Which one matters depends on the question: the
ATT for judging a voluntary programme as it runs, the ATE for making it compulsory.

\section{Selection bias}\label{sec:b2:causality-and-experiments:selection}

\begin{definition}[Selection bias]\label{def:b2:causality-and-experiments:selection-bias}
\emph{Selection bias}\index{selection bias} is the difference $E[Y_i(0) \mid D_i = 1] - E[Y_i(0) \mid D_i = 0]$ between
the untreated outcomes of the treated and of the untreated.
\end{definition}

\begin{proposition}[The selection-bias decomposition]\label{prop:b2:causality-and-experiments:decomposition}
The difference in mean observed outcomes between treated and untreated units equals the ATT plus \omterm{def:b2:causality-and-experiments:selection-bias}{selection bias}:
\[
E[Y_i \mid D_i = 1] - E[Y_i \mid D_i = 0] = \mathit{ATT} + \bigl(E[Y_i(0) \mid D_i = 1] - E[Y_i(0) \mid D_i = 0]\bigr).
\]
\end{proposition}

\begin{proof}
For the treated $Y_i = Y_i(1)$ and for the untreated $Y_i = Y_i(0)$, so the left side is $E[Y_i(1) \mid D_i = 1] -
E[Y_i(0) \mid D_i = 0]$. Add and subtract $E[Y_i(0) \mid D_i = 1]$: the first two terms give $E[Y_i(1) - Y_i(0) \mid D_i =
1] = \mathit{ATT}$, the last two the \omterm{def:b2:causality-and-experiments:selection-bias}{selection bias}.
\end{proof}

The identity also holds exactly in any sample, with sample means in place of expectations. The listing builds both
\omterm{def:b2:causality-and-experiments:potential-outcome}{potential outcomes} for 10,000 simulated workers, lets the less skilled take the training, and checks it.

\omcode{code/bachelor-2/33-causality-and-experiments/listings/selection_sim.py}{1}{26}{Simulated \omterm{def:b2:causality-and-experiments:potential-outcome}{potential outcomes}:
self-selection into training, then random assignment.}\label{lst:b2:causality-and-experiments:selection}

\omoutput{code/bachelor-2/33-causality-and-experiments/outputs/selection_sim.txt}{Training raises earnings, yet the trained
earn less: \omterm{def:b2:causality-and-experiments:selection-bias}{selection bias} outweighs the effect. Random assignment recovers it.}\label{out:b2:causality-and-experiments:selection}

% code: listings/selection_sim.py
The training raises earnings by 1.001 on average, and by 1.284 among those who take it, since its effect is larger for the
less skilled. But those who take it would have earned 2.221 less without it, so the trained earn 0.937 less than the
others. Assigning the same training at random gives a difference of 1.051, close to the ATE.

\section{What randomisation buys}\label{sec:b2:causality-and-experiments:rct}

\begin{definition}[Randomised controlled trial]\label{def:b2:causality-and-experiments:rct}
A \emph{randomised controlled trial}\index{randomised controlled trial} assigns the treatment by a random device, so that
$D_i$ is independent of the \omterm{def:b2:causality-and-experiments:potential-outcome}{potential outcomes} $(Y_i(0), Y_i(1))$.
\end{definition}

With independence, $E[Y_i(0) \mid D_i = 1] = E[Y_i(0) \mid D_i = 0]$: the \omterm{def:b2:causality-and-experiments:selection-bias}{selection bias} is zero, and the ATT equals the
ATE. The difference in means is then an unbiased estimate of the ATE, and its standard error is that of a difference of
two independent means.

\begin{record}[The National Supported Work demonstration]
% ledger: b1:lalonde1986, b2:nsw-cps
In the 1970s an American employment and training programme was run as a field experiment: applicants were assigned at
random to the programme or to a control group. LaLonde (1986) took the men in the experiment and asked whether the
non-experimental methods then common in programme evaluation, applied to the treated men and to comparison samples drawn
from two general surveys, would reproduce the experimental estimate. Many did not. The listings use the subsample of 185
treated and 260 control men distributed by Dehejia and Wahba (1999).
\end{record}

\omcode{code/bachelor-2/33-causality-and-experiments/listings/nsw_experiment.py}{1}{16}{The experimental estimate: a
difference in means, and the same by regression.}\label{lst:b2:causality-and-experiments:nsw}

\omoutput{code/bachelor-2/33-causality-and-experiments/outputs/nsw_experiment.txt}{An effect of 1,794 dollars, with an
interval from 479 to 3,110.}\label{out:b2:causality-and-experiments:nsw}

\begin{definition}[Covariate balance]\label{def:b2:causality-and-experiments:balance}
Treated and control groups show \emph{covariate balance}\index{covariate balance} when the distributions of variables
fixed before treatment are the same in both; randomisation makes them equal in expectation.
\end{definition}

A balance table compares the groups on what was measured before assignment. It cannot check what was not measured, but it
can catch a randomisation that went wrong.

\omcode{code/bachelor-2/33-causality-and-experiments/listings/balance.py}{1}{13}{Pre-treatment means by group, with the
$t$ statistic of each difference.}\label{lst:b2:causality-and-experiments:balance}

\omoutput{code/bachelor-2/33-causality-and-experiments/outputs/balance.txt}{Seven of eight differences are small; the share
without a high-school diploma differs.}\label{out:b2:causality-and-experiments:balance}

% code: listings/balance.py
One difference of eight is significant at 5 per cent: 71 per cent of the trained men have no diploma, against 83 per cent
of the controls. With eight independent comparisons, the chance of at least one such difference is $1 - 0.95^8 = 0.34$
even when the randomisation is perfect.

\section{Regression as adjustment}\label{sec:b2:causality-and-experiments:adjustment}

\begin{definition}[Conditional independence assumption]\label{def:b2:causality-and-experiments:cia}
The \emph{conditional independence assumption}\index{conditional independence assumption} states that, among units with
the same observed characteristics $X_i$, treatment is independent of the \omterm{def:b2:causality-and-experiments:potential-outcome}{potential outcomes}: $(Y_i(0), Y_i(1))$ is
independent of $D_i$ given $X_i$.
\end{definition}

In an experiment, adding pre-treatment controls to the regression changes the estimate only by chance and can make it
more precise. In an observational comparison, controls are the only defence, and they work only if the \omterm{def:b2:causality-and-experiments:cia}{conditional
independence assumption} holds: if, given the controls, nothing that affects earnings also affects who is treated.

\omcode{code/bachelor-2/33-causality-and-experiments/listings/naive_cps.py}{1}{16}{The trained men against men from a
general survey.}\label{lst:b2:causality-and-experiments:naive}

\omoutput{code/bachelor-2/33-causality-and-experiments/outputs/naive_cps.txt}{The survey men are older, better educated,
more often married, and earned nine times as much in 1975.}\label{out:b2:causality-and-experiments:naive}

\omcode{code/bachelor-2/33-causality-and-experiments/listings/regression_adjust.py}{1}{17}{The same controls, in the
experiment and in the survey comparison.}\label{lst:b2:causality-and-experiments:adjust}

\omoutput{code/bachelor-2/33-causality-and-experiments/outputs/regression_adjust.txt}{In the experiment the controls barely
move the estimate; in the survey comparison they move it by 9,291 dollars and still miss.}\label{out:b2:causality-and-experiments:adjust}

% code: listings/naive_cps.py, listings/regression_adjust.py
In the experiment the estimate goes from 1,794 to 1,676 dollars with controls. In the survey comparison it goes from
$-8{,}498$ to 794: the controls remove most of the gap, and the estimate is still less than half the experimental one,
with an interval (from about $-420$ to $2{,}000$) wide enough to include zero. The men who enter a training programme
differ from the general population in ways that age, schooling and past earnings capture only in part. Matching on these
variables, and the other methods that refine this comparison, are the subject of the Year 3 volume.

\section{Classic experiments}\label{sec:b2:causality-and-experiments:classics}

\begin{record}[Tennessee's class-size experiment]
% ledger: b2:krueger1999
In Tennessee's Project STAR, about 11,600 pupils and their teachers were assigned at random to classes of different sizes
from kindergarten to third grade. Krueger (1999), adjusting for pupils who left the experiment or changed class, found that
test scores rose by about four percentile points in the first year in a small class, and that the advantage grew by about
one percentile point a year after that. Teacher aides and measured teacher characteristics had little effect; the effect
of small classes was larger for minority pupils and those eligible for free lunch; and the gains were unlikely to come
merely from being observed.
\end{record}

\section{Whether a result travels}\label{sec:b2:causality-and-experiments:external}

\begin{definition}[External validity]\label{def:b2:causality-and-experiments:external-validity}
A result has \emph{external validity}\index{external validity} if it holds for populations, places, times or versions of
the treatment other than those of the study; \emph{internal validity} is whether it holds for the study's own population.
\end{definition}

An experiment secures internal validity. It says nothing by itself about other places: the men of the NSW sample were a
particular group in a particular decade, and the effect of training on them need not be the effect on others. Where the
experiment is done is itself chosen, and the choice can be related to the result.

\begin{record}[Results that did not travel]
% ledger: b2:allcott2015
Allcott (2015) studied 111 randomised trials of the same programme, letters comparing a household's energy use with its
neighbours', run by American utilities with 8.6 million households in all. Predictions based on the first 10 trials
substantially overstated the programme's effect in the next 101. The first utilities to adopt it served more
environmentally minded areas, whose customers responded more, and utilities first targeted their heaviest users. Each
trial was internally valid; the sequence was not a random sample of sites.
\end{record}

\section{The idea of a natural experiment}\label{sec:b2:causality-and-experiments:natural}

\begin{definition}[Natural experiment]\label{def:b2:causality-and-experiments:natural-experiment}
A \emph{natural experiment}\index{natural experiment} is a situation in which a rule, an accident or a policy, not designed
for research, assigns a treatment in a way that is as good as random for some group, so that comparisons within that group
estimate a causal effect.
\end{definition}

Many questions cannot be settled by an experiment: nobody assigns schooling, minimum wages or interest rates by lottery.
The Year 3 volume studies the designs that use \omterm{def:b2:causality-and-experiments:natural-experiment}{natural experiments}: instruments that shift the treatment without
affecting the outcome directly, comparisons of changes over time between groups exposed and not exposed to a policy, and
comparisons of units just above and just below a threshold that decides treatment. Each replaces randomisation by an
assumption, and each assumption can be argued about with evidence.

\begin{model}{Potential-outcomes model}
\assumes Each unit has two \omterm{def:b2:causality-and-experiments:potential-outcome}{potential outcomes}; only one is observed; a unit's outcomes do not depend on others'
treatment.
\predicts A difference in means equals the effect on the treated plus \omterm{def:b2:causality-and-experiments:selection-bias}{selection bias}; under random assignment the
\omterm{def:b2:causality-and-experiments:selection-bias}{selection bias} is zero.
\tested Experimental and non-experimental estimates on the same treated group can be compared: LaLonde (1986) found that
many non-experimental estimates of the NSW programme's effect missed the experimental one; here a regression on survey
men gives 794 dollars against 1,794.
\breaks Interference between units (displacement in a job market) makes effects depend on the scale of treatment; and an
experiment's effect need not carry to other sites, as Allcott (2015) found.
\end{model}

\section{Exercises}\label{sec:b2:causality-and-experiments:exercises}

\begin{exercise}[$\star$]\label{exo:b2:causality-and-experiments:1}
Read the opening chart. Give each group's mean and interval, and the difference. Do the intervals overlap?
\end{exercise}

\begin{exercise}[$\star$]\label{exo:b2:causality-and-experiments:2}
For a hospital, give an example in which the treated have worse untreated outcomes than the untreated. What sign does the
\omterm{def:b2:causality-and-experiments:selection-bias}{selection bias} have, and what does a raw comparison suggest about the hospital?
\end{exercise}

\begin{exercise}[$\star$]\label{exo:b2:causality-and-experiments:3}
In the balance table, the share without a diploma differs significantly between the groups. Does this show that the
randomisation failed? What would you do?
\end{exercise}

\begin{exercise}[$\star$]\label{exo:b2:causality-and-experiments:4}
A training programme in one town reduces unemployment among participants. Give a reason why making it compulsory in the
whole country might have a smaller effect per participant.
\end{exercise}

\begin{exercise}[$\star\star$]\label{exo:b2:causality-and-experiments:5}
From the listing's output, compute the 95 per cent interval of the experimental estimate from the difference and its
standard error. Is a zero effect inside it?
\end{exercise}

\begin{exercise}[$\star\star$]\label{exo:b2:causality-and-experiments:6}
In the selection listing, which number is the ATE, which the ATT, and why do they differ? Predict what happens to the
\omterm{def:b2:causality-and-experiments:selection-bias}{selection bias} if the more skilled take the training instead.
\end{exercise}

\begin{exercise}[$\star\star$]\label{exo:b2:causality-and-experiments:7}
\textbf{(D)} With the NSW and CPS samples, repeat the survey comparison restricting the survey men to those with 1975
earnings below 5,000 dollars. What happens to the raw and the adjusted estimates? Why?
\end{exercise}

\begin{exercise}[$\star\star$]\label{exo:b2:causality-and-experiments:8}
Why do controls in the experiment move the estimate only from 1,794 to 1,676, while in the survey comparison they move it
by more than 9,000?
\end{exercise}

\begin{exercise}[$\star\star$]\label{exo:b2:causality-and-experiments:9}
A school district reports that pupils who chose a new reading programme scored higher than those who did not. List three
reasons why this might not be the programme's effect, and design an experiment that would settle it.
\end{exercise}

\begin{exercise}[$\star\star\star$]\label{exo:b2:causality-and-experiments:10}
Show that under random assignment the difference in sample means is an unbiased estimator of the ATE, and derive its
variance.
\end{exercise}

\begin{exercise}[$\star\star\star$]\label{exo:b2:causality-and-experiments:11}
Write the difference in means as ATE plus \omterm{def:b2:causality-and-experiments:selection-bias}{selection bias} plus a third term, and interpret the third term.
\end{exercise}

\begin{exercise}[$\star\star\star$]\label{exo:b2:causality-and-experiments:12}
Under the \omterm{def:b2:causality-and-experiments:cia}{conditional independence assumption} with a discrete $X$, show that the ATE equals $\sum_xP(X_i =
x)\bigl(E[Y_i \mid D_i = 1, X_i = x] - E[Y_i \mid D_i = 0, X_i = x]\bigr)$.
\end{exercise}

\section{Problem: Brenland's training-voucher pilot}\label{sec:b2:causality-and-experiments:problem}

\begin{problem}[{Weekend problem --- Brenland's training-voucher pilot: potential outcomes, the estimate and its interval,
a non-random roll-out, and the sample size for a given power}]
\label{pb:b2:causality-and-experiments:1}
Brenland's labour ministry gives training vouchers to workers in a pilot of 800. In the simulation that stands for the
pilot (seeded, in the chapter's code), each worker has an unobserved drive $m_i$, standard normal; earnings without a
voucher are $Y_i(0) = 24{,}000 + 3{,}000m_i + 4{,}000\varepsilon_i$ dalens a year, with $\varepsilon_i$ standard normal; and
the voucher raises earnings by $1{,}500 + 500m_i$.

\medskip
\noindent\textbf{Part I --- Potential outcomes.}
\begin{enumerate}
  \item Write $Y_i(1)$.
  \item What is the ATE in the population model?
  \item Which workers gain most from a voucher?
  \item What is the ATE among the 800 simulated workers?
  \item Why can the ministry never observe a worker's own effect?
  \item Give an example of how the \omterm{def:b2:causality-and-experiments:sutva}{stable unit treatment value assumption} could fail here.
\end{enumerate}

\medskip
\noindent\textbf{Part II --- The estimate and its interval.} Half of the 800, drawn at random, get a voucher.
\begin{enumerate}[resume]
  \item Give the difference in mean earnings.
  \item Give its standard error and 95 per cent interval.
  \item Does the interval contain the true ATE?
  \item Why is the estimate not exactly the ATE?
  \item Why is the ATT equal to the ATE in expectation here?
  \item What would adding pre-treatment earnings as a control do?
\end{enumerate}

\medskip
\noindent\textbf{Part III --- A non-random roll-out.} Instead, workers apply, and those with $m_i + \eta_i > 0$, $\eta_i$
standard normal, get a voucher.
\begin{enumerate}[resume]
  \item What share applies?
  \item Give the difference in mean earnings between voucher holders and others.
  \item Give the ATT among applicants.
  \item Give the \omterm{def:b2:causality-and-experiments:selection-bias}{selection bias}, and check the decomposition.
  \item Why is the ATT larger than the ATE?
  \item Could controlling for education remove the bias? Under what condition?
\end{enumerate}

\medskip
\noindent\textbf{Part IV --- The sample size for a given power.}
\begin{enumerate}[resume]
  \item The ministry wants to detect an effect of 1,500 dalens with 80 per cent power at the 5 per cent level. Write the
  number per arm, $n = 2(z_{0.025} + z_{0.2})^2\sigma^2/\delta^2$.
  \item With $\sigma = 5{,}000$, how many workers per arm?
  \item With the standard deviation of control earnings in the pilot, how many?
  \item To detect 750 dalens, with $\sigma = 5{,}000$, how many?
  \item Why does halving the effect quadruple the sample?
  \item Was the pilot of 800 large enough?
  \item State the result: the ATE estimate with its interval, the naive bias of the roll-out, and the sample size for 80 per
  cent power.
\end{enumerate}
\end{problem}
